% BANKS-SOURCE: pdf=252 print=242 kind=chapter id=chapter11
\BanksChapter{11}{Concluding remarks}

% BANKS-SOURCE: pdf=252 print=242 kind=prose
If you've successfully worked your way through this short but arduous journey to the
world of quantum field theory, you should be exhilarated! The landscape is full of
fabulous beasts (most of which seem to go by the last name -on) and elegant formulae.
Most surprisingly, all of this elegance seems to give us a remarkably precise and general
description of many facets of our world, from phase transitions in humdrum materials
to the interiors of stars and the intergalactic medium. In this concluding section I
want to emphasize again a few general lessons, and chart out for you the parts of the
quantum-field-theory landscape we have NOT explored.

The first of the important lessons that you should take away from this book is the
beautiful unification of the classical theories of fields and particles that is forced on
us by combining relativity and quantum mechanics in a fixed space-time background.
The second is the unification of the methods of quantum field theory and classical
statistical mechanics, which is provided by the Euclidean path-integral formulation of
field theory.

Next I would ask you to remember the difference between a symmetry and a gauge
equivalence and the different meanings of the idea of spontaneous symmetry breakdown
in the two cases. Spontaneous breakdown of a global symmetry is related to locality. A
quantum field theory is defined by its behavior at short distances, but there may be different
infrared realizations of the same short-distance operator algebra and Hamiltonian.
Sometimes these are related by a symmetry transformation. If it is a continuous group
of transformations, any one infrared sector of the theory contains massless excitations
called Nambu--Goldstone particles.

By contrast, no physical quantity transforms under a gauge equivalence, which is
merely a convenient redundancy of our description of a physical system. What we call
a spontaneously broken gauge symmetry is a particular phase, the Higgs phase, of gauge
theories. In this phase there are no massless particles associated with the generators
of the gauge equivalence. Gauge theories can also have another massive phase (which
is distinct from the Higgs phase if the Higgs fields are invariant under a non-trivial
subgroup of the center of the gauge group) called the confining phase, which is related
to the Higgs phase by electric--magnetic duality. In the real world the non-abelian
electro-weak gauge group is in the Higgs phase, and the color group of QCD is in its
confining phase. The U(1) group of electromagnetism is in yet a third phase, called the
Coulomb phase. This is a special case of a more general possibility in which the gauge
theory behaves like a conformal field theory in the infrared.

% BANKS-SOURCE: pdf=253 print=243 kind=prose
Perhaps the most important topic in quantum field theory is the theory of renormalization.
As formulated by Wilson, Kadanoff, and Fisher, this is a very general way
of understanding how physical processes at different length scales are related to each
other. From the philosophical point of view, this is the reason why we are able to do
physics at all. If we had to uncover the correct theory of the smallest length scales and
highest energies in order to talk about what we see in our laboratories, there would be
no hope for a theory of physics at all. The renormalization group (RG) explains
to us why it is that we can do long-distance physics without knowing about short distances.
The key concept is that of the RG fixed point, parametrized by a small number
of relevant and marginal perturbations. All UV theories containing a given IR set of
degrees of freedom flow to the vicinity of the fixed point whose basin of attraction they
lie in. They are differentiated only by the values they determine for the marginal and
relevant parameters. This gives rise to universality classes of behavior, a post-diction
spectacularly confirmed by observations of critical phenomena in condensed-matter
physics.

Given this broad-brush picture of what we have covered in this book, we now turn
to the subjects we have omitted. The most important of these are supersymmetry,
finite-temperature field theory, and field theory in curved space-time. Supersymmetry
is (in dimensions higher than 2) the only allowed extension of the space-time conformal
group. Remarkably, it joins together bosons and fermions as a single entity. It also seems
to be deeply connected to the quantum theory of gravity. At the technical level, supersymmetry
allows us to make many exact statements about quantum field theory that are
not possible with ordinary field theories. There are several good reviews of SUSY
field theory \cite{banks-ref-152,banks-ref-153,banks-ref-154,banks-ref-155,banks-ref-156}, but, if I may be
allowed an opinion, as yet no comprehensive treatment of all modern developments.

Finite-temperature relativistic field theory is primarily applicable in cosmology and
astrophysics, though the techniques are similar to those of the non-relativistic theory,
which has wide applications in condensed-matter physics. Equilibrium calculations are
quite similar to what we have discussed: to calculate averages in the canonical ensemble,
instead of using the vacuum density matrix, one simply performs the Euclidean
path integral on a space with one compactified dimension, whose length is the inverse
temperature. Good reviews of this subject can be found in \cite{banks-ref-157}. Non-equilibrium field
theory is much more difficult, and has been the focus of a lot of recent work. As far as
I know, there is no definitive modern summary, but the reader can consult \cite{banks-ref-158}.

Quantum field theory in curved space-time is relatively easy to define for space-times
whose complexification has a real Euclidean section. Again, one simply performs the
path integral on the appropriate Euclidean manifold, and analytically continues the
answer. More general space-times with no time-like isometries can also be dealt with
at a certain level of generality \cite{banks-ref-159}. However, the most interesting thing about this
subject is that it contains the seeds of its own demise, and shows us that we need to
replace quantum field theory with a quantum theory of gravity. I am referring to the
phenomenon of Hawking radiation. Within the context of field theory in curved space-time,
it seems to imply that the formation and evaporation of black holes violates the
unitary evolution postulate of quantum mechanics. It is only with the advent of string

% BANKS-SOURCE: pdf=254 print=244 kind=prose
theory, the first real theory of quantum gravity, that we have begun to understand how
this paradox is resolved.

\BanksImplicitHook{I-CH11-001}

Other subjects to which we have given short shrift are perturbative QCD \cite{banks-ref-124}
and weak interactions \cite{banks-ref-160,banks-ref-161}. These subjects are nicely treated in the textbook by
Peskin and Schroeder \cite{banks-ref-33}. Lattice field theory has developed into a computer-intensive
subfield. Good reviews of the subject can be found in \cite{banks-ref-162}. Appendix A contains references
to a number of books on the vast subject of quantum field theory in statistical
physics. Finally, I want to mention the use of two-dimensional conformal field theory
in statistical physics and string theory. This is a subject of great beauty and utility
\cite{banks-ref-163,banks-ref-164}.

I hope you've come to the end of this trip with an appreciation of the beauties of
quantum field theory and a hunger to know more about it. As you can see, there are
lots of directions to follow from this point. The methods of field theory and the concept
of renormalization have become so pervasive that it is probably beyond the capabilities
of any single person to be an expert in the entire field. But every single avenue you
can follow is interesting and exciting. Even those of you who prefer to explore the
mysteries of quantum gravity, where the paradigms of field theory appear to fail, will
continue to use many of the ideas and techniques of field theory. Indeed, we have found
that, in many cases, the dynamics of gravity in some class of space-times is completely
equivalent to that of a quantum field theory living on an auxiliary space
\cite{banks-ref-165,banks-ref-166,banks-ref-167}.
